\documentclass[10pt,a4paper]{amsart}
\usepackage{amsmath,amssymb,amsthm,graphicx}
\usepackage{verbatim}
\newenvironment{reference}{\comment}{\endcomment}
%\newenvironment{reference}{}{}
\newenvironment{slogan}{\comment}{\endcomment}
\newenvironment{history}{\comment}{\endcomment}

% For commutative diagrams we use Xy-pic
\usepackage[all]{xy}

% We use 2cell for 2-commutative diagrams.
\xyoption{2cell}
% Theorem environments.
%
\theoremstyle{plain}
\newtheorem{theorem}[subsection]{Theorem}
\newtheorem{proposition}[subsection]{Proposition}
\newtheorem{lemma}[subsection]{Lemma}

\theoremstyle{definition}
\newtheorem{definition}[subsection]{Definition}
\newtheorem{example}[subsection]{Example}
\newtheorem{exercise}[subsection]{Exercise}
\newtheorem{situation}[subsection]{Situation}

\theoremstyle{remark}
\newtheorem{remark}[subsection]{Remark}
\newtheorem{remarks}[subsection]{Remarks}

\numberwithin{equation}{subsection}

% Macros
%
\def\lim{\mathop{\mathrm{lim}}\nolimits}
\def\colim{\mathop{\mathrm{colim}}\nolimits}
\def\Spec{\mathop{\mathrm{Spec}}}
\def\Hom{\mathop{\mathrm{Hom}}\nolimits}
\def\Ext{\mathop{\mathrm{Ext}}\nolimits}
\def\SheafHom{\mathop{\mathcal{H}\!\mathit{om}}\nolimits}
\def\SheafExt{\mathop{\mathcal{E}\!\mathit{xt}}\nolimits}
\def\Sch{\mathit{Sch}}
\def\Mor{\mathop{\mathrm{Mor}}\nolimits}
\def\Ob{\mathop{\mathrm{Ob}}\nolimits}
\def\Sh{\mathop{\mathit{Sh}}\nolimits}
\def\NL{\mathop{N\!L}\nolimits}
\def\CH{\mathop{\mathrm{CH}}\nolimits}
\def\proetale{{pro\text{-}\acute{e}tale}}
\def\etale{{\acute{e}tale}}
\def\QCoh{\mathit{QCoh}}
\def\Ker{\mathop{\mathrm{Ker}}}
\def\Im{\mathop{\mathrm{Im}}}
\def\Coker{\mathop{\mathrm{Coker}}}
\def\Coim{\mathop{\mathrm{Coim}}}

% Boxtimes
%
\DeclareMathSymbol{\boxtimes}{\mathbin}{AMSa}{"02}

%
% Macros for moduli stacks/spaces
%
\def\QCohstack{\mathcal{QC}\!\mathit{oh}}
\def\Cohstack{\mathcal{C}\!\mathit{oh}}
\def\Spacesstack{\mathcal{S}\!\mathit{paces}}
\def\Quotfunctor{\mathrm{Quot}}
\def\Hilbfunctor{\mathrm{Hilb}}
\def\Curvesstack{\mathcal{C}\!\mathit{urves}}
\def\Polarizedstack{\mathcal{P}\!\mathit{olarized}}
\def\Complexesstack{\mathcal{C}\!\mathit{omplexes}}
% \Pic is the operator that assigns to X its picard group, usage \Pic(X)
% \Picardstack_{X/B} denotes the Picard stack of X over B
% \Picardfunctor_{X/B} denotes the Picard functor of X over B
\def\Pic{\mathop{\mathrm{Pic}}\nolimits}
\def\Picardstack{\mathcal{P}\!\mathit{ic}}
\def\Picardfunctor{\mathrm{Pic}}
\def\Deformationcategory{\mathcal{D}\!\mathit{ef}}
\usepackage{cleveref}
\crefname{originalsection}{original section}{original sections}
\Crefname{originalsection}{Original section}{Original sections}
\makeatletter
\newlabel{constructions-section-grassmannian}{{089R}{source}{}{sourcepointer-1}{}}
\newlabel{constructions-section-grassmannian@cref}{{[originalsection][0][]089R}{[1][0][]source}{}{}{}}
\makeatother
\crefname{originalreference}{original reference}{original references}
\Crefname{originalreference}{Original reference}{Original references}
\makeatletter
\newlabel{lemma-enough}{{00YM}{source}{}{sourcepointer-2}{}}
\newlabel{lemma-enough@cref}{{[originalreference][0][]00YM}{[1][0][]source}{}{}{}}
\makeatother
\crefname{originalreference}{original reference}{original references}
\Crefname{originalreference}{Original reference}{Original references}
\makeatletter
\newlabel{lemma-grassmanian}{{0FIQ}{source}{}{sourcepointer-3}{}}
\newlabel{lemma-grassmanian@cref}{{[originalreference][0][]0FIQ}{[1][0][]source}{}{}{}}
\makeatother
\crefname{originallemma}{original lemma}{original lemmas}
\Crefname{originallemma}{Original lemma}{Original lemmas}
\makeatletter
\newlabel{more-morphisms-lemma-two-unramified-morphisms-formally-smooth}{{067W}{source}{}{sourcepointer-4}{}}
\newlabel{more-morphisms-lemma-two-unramified-morphisms-formally-smooth@cref}{{[originallemma][0][]067W}{[1][0][]source}{}{}{}}
\makeatother
\title{A sufficient condition for compatibility of the cycle class map with pushforward}
\author{The Stacks Project Authors}
\date{stacks-project:weil.tex:local-snapshot-20261012; GNU FDL 1.2 or later}
\begin{document}
\maketitle
\noindent\textit{Editorial note.} All delivered mathematical text is the author's own, copied from the pinned Stacks Project chapter file weil.tex; no statement or proof step was rewritten, reordered or removed.\par\medskip
\noindent\textit{Editorial note.} The counted proof is the author's complete argument: the reduction to a closed immersion, the use of a Grassmannian to convert the conormal class into a pullback class, and the conormal sequence verifying the hypothesis.\par\medskip
\noindent\textit{Editorial note.} The label lemma-enough-better names the counted statement, and the proof environment immediately following it in the source is delivered verbatim as its proof.\par\medskip
\noindent\textit{Editorial note.} The counted statement is carried by the author's own lemma environment in the chapter source.\par\medskip
\noindent\textit{Editorial note.} Pre-open gate: python operations/agent-authoring/tools/precard\_gate.py --source weil.tex --label lemma-enough-better --start 5257 --end 5300 --local-id 28167 returned RESULT=OK. Pre-assembly tag self-check: canonical tag resolved and not present in reference/public-items.jsonl.\par\medskip
\noindent\textit{Editorial note.} The statement and proof are copied from the pinned chapter file; the cited inputs are kept as external pointers with their canonical Stacks tags.\par\medskip
\noindent\textit{Allowed standard prerequisites (editorial).}\begin{itemize}
\item Reduction to closed immersions. Same chapter, the lemma reducing the compatibility of the cycle class map with pushforward to the case of a closed immersion. Reduces the general statement to closed immersions between nonempty smooth projective equidimensional schemes.
\item Compatibility for Grassmannian projections. Same chapter, the lemma that the cycle class map commutes with pushforward along the projection from a product with a Grassmannian. Allows pushing the closed immersion into a product with a Grassmannian without changing the cycle class comparison.
\item Conormal sequence for a closed immersion into a product. More on Morphisms, the lemma describing the conormal sheaf of a closed immersion given by two formally smooth morphisms, and Constructions, the construction of the morphism to the Grassmannian from a surjection onto a twisted conormal sheaf. Shows that after twisting, the conormal sheaf becomes the pullback of a class, which is the hypothesis of the reduction.
\end{itemize}\medskip
\ifdefined\hypertarget\hypertarget{sourcepointer-1}{}\fi
\noindent\textit{Source reference.} Original section 089R: sources/stacks-chow/tags.txt, source lines 10506--10506. Prerequisite pointer invoked at proof line(s) 5279; it names a standard fact of a sibling chapter that the counted argument uses, and every step of the counted proof built on it is retained verbatim in the delivered body.\par\medskip
\ifdefined\hypertarget\hypertarget{sourcepointer-2}{}\fi
\noindent\textit{Source reference.} Original reference 00YM: sources/stacks-weil/pinned-original-source.tex, source lines 5103--5103. Prerequisite pointer invoked at proof line(s) 5268; it names a standard fact of the same chapter that the counted argument uses, and every step of the counted proof built on it is retained verbatim in the delivered body.\par\medskip
\ifdefined\hypertarget\hypertarget{sourcepointer-3}{}\fi
\noindent\textit{Source reference.} Original reference 0FIQ: sources/stacks-weil/pinned-original-source.tex, source lines 5228--5228. Prerequisite pointer invoked at proof line(s) 5285; it names a standard fact of the same chapter that the counted argument uses, and every step of the counted proof built on it is retained verbatim in the delivered body.\par\medskip
\ifdefined\hypertarget\hypertarget{sourcepointer-4}{}\fi
\noindent\textit{Source reference.} Original lemma 067W: sources/stacks-chow/tags.txt, source lines 7985--7985. Prerequisite pointer invoked at proof line(s) 5293; it names a standard fact of a sibling chapter that the counted argument uses, and every step of the counted proof built on it is retained verbatim in the delivered body.\par\medskip
\begin{lemma}
\label{lemma-enough-better}
Assume given data (D0), (D1), and (D2') satisfying axioms (A1) -- (A7).
In order to show that $\gamma$ commutes with pushforward it suffices
to show that $i_*(1) = \gamma([Z])$ if $i : Z \to X$ is a closed
immersion of nonempty smooth projective equidimensional schemes over $k$
such that the class of $\mathcal{C}_{Z/X}$ in $K_0(Z)$ is the
pullback of a class in $K_0(X)$.
\end{lemma}

\begin{proof}
By Lemma \ref{lemma-enough} it suffices to show that $i_*(1) = \gamma([Z])$
if $i : Z \to X$ is a closed immersion of nonempty
smooth projective equidimensional
schemes over $k$. Say $Z$ has codimension $r$ in $X$.
Let $\mathcal{L}$ be a sufficiently ample invertible module on $X$.
Choose $n > 0$ and a surjection
$$
\mathcal{O}_Z^{\oplus n} \to \mathcal{C}_{Z/X} \otimes \mathcal{L}|_Z
$$
This gives a morphism $g : Z \to \mathbf{G}(n - r, n)$
to the Grassmannian over $k$, see
Constructions, Section \ref{constructions-section-grassmannian}.
Consider the composition
$$
Z \to X \times \mathbf{G}(n - r, n) \to X
$$
Pushforward along the second morphism is compatible with classes
of cycles by Lemma \ref{lemma-grassmanian}. The conormal sheaf $\mathcal{C}$
of the closed immersion $Z \to X \times \mathbf{G}(n - r, n)$ sits in
a short exact sequence
$$
0 \to \mathcal{C}_{Z/X} \to \mathcal{C} \to
g^*\Omega_{\mathbf{G}(n - r, n)} \to 0
$$
See More on Morphisms, Lemma
\ref{more-morphisms-lemma-two-unramified-morphisms-formally-smooth}.
Since $\mathcal{C}_{Z/X} \otimes \mathcal{L}|_Z$ is the pull
back of a finite locally free sheaf on $\mathbf{G}(n - r, n)$
we conclude that the class of $\mathcal{C}$ in $K_0(Z)$
is the pullback of a class in $K_0(X \times \mathbf{G}(n - r, n))$.
Hence we have the property for $Z \to X \times \mathbf{G}(n - r, n)$
and we conclude.
\end{proof}

\end{document}
